\chapter{Resultados e Discussão}
\label{cap4_resultados}

% ---------------------------------------------------------------
% ORIENTAÇÃO: se o seu trabalho tiver um capítulo intermediário
% (por exemplo, "Desenvolvimento do sistema" ou "O caso"), crie um
% novo arquivo em Capitulos/ e inclua-o no main.tex antes deste.
% Se os resultados forem extensos, a discussão pode virar um
% capítulo próprio.
% ---------------------------------------------------------------

Este capítulo apresenta os resultados obtidos e os discute à luz dos objetivos e da literatura.

\section{Apresentação dos resultados}
\label{sec_res_apresentacao}

Apresente os resultados de forma objetiva, organizados por objetivo específico ou por questão de pesquisa. Use figuras, quadros e tabelas sempre que facilitarem a compreensão, e cite cada um no texto antes de ele aparecer, comentando o que deve ser observado. Os exemplos do Capítulo~\ref{cap2_fundamentacao} mostram formas usuais de apresentação para cada tipo de trabalho, como a Tabela~\ref{tab_resultados_experimento} e a Figura~\ref{fig_likert}.

\section{Discussão}
\label{sec_res_discussao}

Interprete os resultados: o que eles significam, se respondem à questão de pesquisa, se confirmam ou contrariam os trabalhos relacionados (Seção~\ref{sec_trabalhos_relacionados}) e quais as suas implicações práticas. A discussão é o que diferencia uma monografia de um relatório técnico.

\section{Ameaças à validade}
\label{sec_res_ameacas}

Descreva os fatores que podem ter influenciado os resultados e o que foi feito para reduzi-los, por exemplo: tamanho e representatividade da amostra, viés do pesquisador, configuração do ambiente de testes e possibilidade de generalização dos resultados \cite{wohlin2012experimentation}.
